\subsection{Privilege Separation、Monitoring、DataSentinel 与 Provenance}

slides 86--98 收束剩余防线。身份与访问控制限制谁可触发任务；task decomposition 把高权限步骤拆给隔离组件；Privtrans 自动划分程序实现 privilege separation；runtime monitoring 观察异常行动与信息流；DataSentinel 用博弈论检测 prompt injection；information-flow control 限制敏感数据传播；provenance 记录数据和决策来源，支持审计与恢复。最后一页总结攻击、评测和防御必须共同演进。

\lecturefigure{slides-images/slide-086.png}{官方 slide 86：Defense mechanisms 进入身份与访问控制}
\lecturefigure{slides-images/slide-087.png}{官方 slide 87：按 identity 与 policy 管理 user access}
\lecturefigure{slides-images/slide-089.png}{官方 slide 89：Task decomposition 与 privilege separation}
\lecturefigure{slides-images/slide-090.png}{官方 slide 90：Privtrans 自动权限分离}
\lecturefigure{slides-images/slide-092.png}{官方 slide 92：Runtime monitoring 与异常行为观察}
\lecturefigure{slides-images/slide-093.png}{官方 slide 93：DataSentinel 博弈论 Prompt Injection 检测}
\lecturefigure{slides-images/slide-095.png}{官方 slide 95：Information flow monitoring/control}
\lecturefigure{slides-images/slide-097.png}{官方 slide 97：构建 provenance 的防御位置}
\lecturefigure{slides-images/slide-098.png}{官方 slide 98：Agentic AI 安全与安全性的课程结论}

\begin{knowledgebox}{读图：最后四层解决不同时间尺度}
Policy enforcement 在动作前阻断，privilege separation 限制单次失守的爆炸半径，runtime monitoring 在执行中检测偏离，provenance 在事后支持归因、撤销和训练回流。只做其中一层会留下时间窗口：例如仅靠监控发现问题时，危险写操作可能已完成。
\end{knowledgebox}

\begin{importantbox}{最小可信执行链}
每个高风险 action 都应能回答：谁提出、基于哪些来源、经过哪条 policy、使用何种权限、修改了什么状态、结果如何验证、怎样撤销。若任一项不可回答，系统就还没有达到可审计部署标准。
\end{importantbox}

\begin{warningbox}{老师强调：安全是持续过程}
模型、工具、外部网站和攻击方法都在变化，一次通过 red team 不代表长期安全。攻击轨迹必须进入回归集，policy 与 privilege 变化要版本化，监控告警要有值班与修复 SLA，provenance 要能支撑事故复盘和用户补偿。
\end{warningbox}

